\section{Results}\label{sec:results}

\subsection{A catalogue the declared taxonomy does not supply (RQ1)}\label{sec:res-rq1}

Clustering the 5{,}243 cases yields between 40 and 120 types across the sixteen configurations,
with no singleton clusters under any of them. On the raw arm, density-based clustering produces
64 types with a median size of 50 cases and 879 cases left unassigned; centroid clustering
produces 64 types with a median size of 70 and no unassigned cases; consensus produces 120 types
with a median size of 33. Cluster-wise stability, assessed by drawing 25 subsamples of 80\% of
the cases, refitting, and matching each original cluster to its best Jaccard overlap
\citep{hennig2007stability}, gives a mean recovery of 0.7526 on the raw arm and 0.7724 on the
anonymised arm, with 47 of 64 and 47 of 61 types recovered above 0.60. Two limits attach to that
figure. It is a lower bound, because the clusterer is refit on each subsample but the
dimensionality reduction is not; and it was computed for the two density-based clusterings only,
so the remaining fourteen carry no stability estimate.

Against the administration's own taxonomy, the emergent partitions agree moderately, and the
level depends on which clustering is used. Across the eight main clusterings, ARI ranges from
0.086 to 0.202 and purity from 0.732 to 0.759. The majority-class baseline differs by scope and
both are reported: 0.302 for the six clusterings scored on all 5{,}243 cases, and 0.276 for the
two density-based ones, which are scored only on the cases they assign (4{,}364 and 4{,}504 of
5{,}243, the remainder being unassigned as noise). The upper end of the ARI range belongs to the
density-based clusterings; the algorithms that assign every case reach ARI 0.086 to 0.129. The
population confound identified in Section~\ref{sec:res-rq2} applies here as well, and we report
both rather than quoting the flattered figure alone.

The result that does not depend on that choice is what happens when one label is removed. A
single catch-all --- covering official documents of every kind, and applied to 30.5\% of the
5{,}445 concluded cases and 30.2\% of the 5{,}243 analysed --- accounts for most of the
disagreement. Excluding it, ARI rises to 0.352 for the density clusterings and to 0.166--0.256
for the full-coverage ones: a relative increase of 74\% to 110\%, in every case.

That increase needs a null before it can carry the reading we give it, because removing any
large class changes the index. Table~\ref{tab:catch-all-null} supplies one: for each clustering
we drop a randomly chosen set of whole declared classes totalling within 10\% of the catch-all's
size, and rescore. Removing random classes of that size \emph{lowers} agreement --- null means
of 0.064 to 0.160 against unrestricted ARI of 0.086 to 0.202 --- while removing the catch-all
raises it to 0.166--0.352. Across eight clusterings and roughly seven hundred draws each,
\textbf{no null draw reaches the observed value}, at $z$ between 7.3 and 17.2.

That null has a limitation we should name rather than let a reader discover. The catch-all holds
1{,}663 cases and no other declared class exceeds 764, so a size-matched draw is necessarily a
\emph{union} of many classes --- a median of twenty-seven. The comparison it makes is therefore
between removing one large heterogeneous class and removing many small ones, which is not
cleanly a test of \emph{which} class was removed.

The contrast that does isolate that is simpler, and it is unanimous. Dropping each of the three
next-largest declared classes on its own \emph{lowers} agreement in every case --- 24 of 24
single-class drops across the eight clusterings --- while dropping the catch-all raises it in
8 of 8. On the raw-arm centroid clustering, agreement is 0.128 unrestricted, 0.246 without the
catch-all, and 0.107, 0.122 and 0.118 without each of the next three classes in turn. Removing a
large class is not what produces the effect, and removing a large class is in fact mildly
harmful; removing \emph{this} class is what produces it.

The declared taxonomy is therefore not uninformative. It degrades concentrated in one class,
which is what a filing classification does when it is maintained for routing paperwork rather
than for describing procedure. The reading that class is a defect of the taxonomy is ours and
should be resisted: a classification maintained to route documents is fit for the purpose it
serves, and we are scoring it against an analytical purpose it was never designed for.

%% GENERATED by `pipeline run results` - do not edit by hand.
\begin{table}[htbp]
\centering
\caption{The catch-all class against a same-size random-class null.}\label{tab:catch-all-null}
\begin{tabular}{lrrrrrr}
\toprule
clustering & ARI all & ARI excl. & null mean & null max & $\geq$ obs. & $z$ \\
\midrule
exp-raw/hdbscan & 0.1383 & 0.2767 & 0.1010 $\pm$ 0.0189 & 0.1311 & 0/724 & 9.3 \\
exp-raw/kmeans & 0.1324 & 0.2548 & 0.0988 $\pm$ 0.0106 & 0.1231 & 0/724 & 14.8 \\
exp-raw/leiden & 0.1197 & 0.2348 & 0.0885 $\pm$ 0.0102 & 0.1076 & 0/724 & 14.3 \\
exp-anon/hdbscan & 0.1450 & 0.2831 & 0.1072 $\pm$ 0.0187 & 0.1375 & 0/724 & 9.4 \\
exp-anon/kmeans & 0.1317 & 0.2585 & 0.0974 $\pm$ 0.0110 & 0.1225 & 0/724 & 14.6 \\
exp-anon/leiden & 0.1161 & 0.2303 & 0.0851 $\pm$ 0.0102 & 0.1068 & 0/724 & 14.2 \\
\bottomrule
\end{tabular}
\par\vspace{2pt}\footnotesize\raggedright Null: drop a randomly chosen set of whole declared classes totalling within 10\% of the catch-all's size, and rescore. Removing random classes of the same size lowers agreement; removing the catch-all raises it. No null draw reaches the observed value for any clustering.
\end{table}


\subsubsection{What an entry looks like}

An artefact offered as a transparency instrument should be shown, not only measured.
Table~\ref{tab:catalogue} gives ten entries from the partition Section~\ref{sec:selection}
selects as the deliverable --- the six largest evaluated types and the four smallest --- with
the generated name, the number of cases assigned, and the support measures for that type's own
typical path. Names are reproduced as the pipeline produced them, and are unvalidated
(Section~\ref{sec:naming}).

The table is included because the averages elsewhere in this section conceal how unevenly
evidenced the entries are, and a reader deciding whether to trust the artefact needs to see the
spread rather than the mean. The best entries are very good: a company-registration type over
325 cases reaches 0.974 next-step accuracy with complete unit recall over a four-visit path and
no fabricated transitions. Others are weak in ways a published catalogue would have to disclose.
A leave-and-indemnity type over 261 cases reaches 0.522 with unit recall 0.515, and a ponto-justification type over 137 cases reaches 0.424 --- both summarising in three or four
visits a set of cases that evidently do more than that. A property-rental payment type covers
296 cases with a 28-visit path that recovers 0.917 of what its cases do and is the only entry in
the sample that fabricates at all.

Two things follow for anyone publishing such a catalogue. An entry is not a uniform object, so
case counts and support measures belong beside every published type rather than in an appendix.
And the weak entries are not the small ones --- three of the four weakest here have more than
130 cases each --- so size alone is not a usable proxy for whether an entry can be trusted.

%% GENERATED by `pipeline run results` - do not edit by hand.
\begin{table}[htbp]
\centering
\caption{Illustrative catalogue entries, largest and smallest evaluated types.}\label{tab:catalogue}
\footnotesize
\begin{tabular}{p{0.42\textwidth}rrrr}
\toprule
process type (generated name) & cases & eval. & top-1 & unit rec. \\
\midrule
Gestão de Contracheques e Solicitações de Pagamentos Indenizatórios & 318 & 312 & 0.850 & 0.978 \\
Solicitação de Autorização para Adesão a Ata de Registro de Preços & 229 & 217 & 0.733 & 0.823 \\
Registro de Alterações Contratuais e Consolidação Societária de Empresas & 210 & 210 & 0.984 & 1.000 \\
Solicitação de Pagamento de Serviços de Mão de Obra Terceirizada & 198 & 198 & 0.771 & 0.924 \\
Retirada de Descontos Indevidos em Contracheques de Servidores & 188 & 188 & 0.770 & 0.916 \\
Gestão de Pessoal e Contratações Temporárias & 192 & 185 & 0.718 & 0.658 \\
Justificativa e Abono de Faltas por Doença ou Trabalho Remoto & 42 & 39 & 0.547 & 0.810 \\
Autorização de Recebimento de Bens Arrematados em Leilão & 38 & 35 & 0.640 & 0.816 \\
Solicitação de Pagamento de Fornecimento de Água Tratada & 31 & 31 & 0.785 & 0.977 \\
Solicitação de Abono de Faltas e Justificativa de Ponto & 30 & 30 & 0.570 & 0.857 \\
\bottomrule
\end{tabular}
\par\vspace{2pt}\footnotesize\raggedright Illustrative entries from the catalogue, largest and smallest evaluated types. Names are generated (Section 4.7) and unvalidated; measures are for the GRU path on the routing alphabet, the family that gains most on routing (Section 4.8).
\end{table}


We do not claim the emergent catalogue is correct. There is no ground truth against which it
could be scored --- the absence of a trustworthy classification is the premise of the work.
What can be established is that it is not arbitrary (Section~\ref{sec:res-rq3}) and that it is
reproducible: all nineteen clusterings defined as a function of the feature matrix refit exactly
from their stamped inputs (the two consensus partitions are derived from other labels, not from
features, and are outside that check).

%% GENERATED by `pipeline run results` - do not edit by hand.
\begin{table}[htbp]
\centering
\caption{Agreement between the emergent partition and the declared taxonomy.}\label{tab:taxonomy}
\footnotesize
\begin{tabular}{llrrrrr}
\toprule
clustering & scope & $n$ & majority & purity & NMI & ARI \\
\midrule
exp-anon/hdbscan & all & 5,243 & 0.302 & 0.730 & 0.530 & 0.145 \\
exp-anon/hdbscan & excl.\ catch-all & 3,660 & 0.208 & 0.775 & 0.621 & 0.283 \\
exp-anon/kmeans & all & 5,243 & 0.302 & 0.735 & 0.531 & 0.132 \\
exp-anon/kmeans & excl.\ catch-all & 3,660 & 0.208 & 0.784 & 0.624 & 0.259 \\
exp-anon/leiden & all & 5,243 & 0.302 & 0.727 & 0.521 & 0.116 \\
exp-anon/leiden & excl.\ catch-all & 3,660 & 0.208 & 0.776 & 0.613 & 0.230 \\
exp-raw/hdbscan & all & 5,243 & 0.302 & 0.728 & 0.527 & 0.138 \\
exp-raw/hdbscan & excl.\ catch-all & 3,660 & 0.208 & 0.775 & 0.621 & 0.277 \\
exp-raw/kmeans & all & 5,243 & 0.302 & 0.738 & 0.531 & 0.132 \\
exp-raw/kmeans & excl.\ catch-all & 3,660 & 0.208 & 0.784 & 0.622 & 0.255 \\
exp-raw/leiden & all & 5,243 & 0.302 & 0.723 & 0.524 & 0.120 \\
exp-raw/leiden & excl.\ catch-all & 3,660 & 0.208 & 0.780 & 0.618 & 0.235 \\
\bottomrule
\end{tabular}
\end{table}


\subsection{Feature blocks and temporal dependency}\label{sec:res-ablation}

Cases are clustered from their complete activity history. This is a design choice rather than an
oversight: the catalogue is a retrospective account, and because the study declines to rely on
the declared type, the partition must use whatever signal the record affords. The ablation in
Table~\ref{tab:ablation} asks a different question --- when does that signal become available?

Removing the keyword block perturbs the partition about as much as discarding half of every
case's history (ARI 0.414 against 0.430), so the controlled vocabulary contributes signal the
other blocks do not carry. Removing the topic block costs less (0.577), which is expected: topic
and semantic representations are both derived from the same opening document.

Truncating each case's features at 75\%, 50\% and 25\% of its own timeline degrades agreement
with the full-history partition smoothly --- ARI 0.581, 0.430 and 0.296 --- rather than
collapsing. A partition built from partial histories is still substantially the same partition,
but a real share of the evidence arrives late. Read as a statement about administrative work
rather than about method, this says that what a case \emph{is} becomes progressively clearer as
it runs: routing history and accumulated actions carry type information that the opening
document does not fully determine.

Two consequences follow, both of them scope statements. The catalogue is retrospective by
construction, which is what an as-is instrument should be. And assigning a type at intake is a
distinct and harder problem: at a 25\% cutoff the partition retains ARI 0.296 against the
full-history one, informative but far from equivalent. Doing that well would require richer
early-stage features than the single opening document embedded here.

\subsubsection{Does the partition predict the future because it has seen it?}\label{sec:holdout}

Clustering on complete histories raises a sharper objection than scope, and it bears directly on
RQ3 rather than on RQ1. Two of the six feature blocks are built from the organisational units a
case visits and the transitions between them, so on the routing alphabet a held-out case is
assigned to a cluster partly by the very sequence the model is then asked to predict. If that is
what produces the routing gain, the gain measures circularity rather than process structure.
Held-out splitting does not address this, because the leak is inside the feature construction
rather than across the split.

The ablation clusterings make the objection partly testable. Each truncated variant builds its
features from the first fraction of every case's own timeline; we then predict on the \emph{full}
held-out traces. The test is weaker than it first appears and we state its limit before its
result: truncating at 25\% still lets the unit blocks see the first quarter of the sequence the
model is scored on from step one, so the leak is reduced rather than removed. What the
comparison can show is whether the gain depends on the \emph{later} part of the history, not
whether it is free of the earlier part.

It does not, and the naive reading of the gradient is wrong in a way worth reporting. Scored as
each truncation was evaluated, the routing gain falls from 0.0381 on full histories to 0.0145 at
25\%. But those four rows are computed on four different populations --- 962, 946, 897 and 889
held-out cases --- and the fall is almost entirely the baseline moving, not the partition losing
information. Per-cluster accuracy is flat across the four truncations (0.6576 to 0.6572, a
change of 0.0004), while the global model \emph{rises} from 0.6195 to 0.6427. A single global
chain cannot change with the clustering, so that rise is the evaluated population becoming
easier as truncation shrinks it.

Rescoring every truncation on the 637 cases all four evaluate --- the same correction this
section applies to the algorithm comparison --- gives gains of 0.0380, 0.0400, 0.0289 and
0.0310. Roughly \textbf{82\% of the routing gain survives when the partition is built from the
first quarter of each case}, and the residual gradient is not monotone: the 25\% truncation
scores above the 50\% one. The action gain stays near zero throughout, consistent with
Section~\ref{sec:res-rq3}.

Three limits bound what this establishes. The truncation is partial, as stated above. The
matched population is the easier subset that survives every truncation, so it is not the corpus.
And the test runs on one arm and one algorithm. Our conclusion is correspondingly narrow: the
routing result is not primarily an artefact of clustering on the later history, and we cannot
show it is free of the early history at all.

%% GENERATED by `pipeline run results` - do not edit by hand.
\begin{table}[htbp]
\centering
\caption{Downstream gain when the partition sees only part of each case's history.}\label{tab:temporal-holdout}
\begin{tabular}{lrrrrr}
\toprule
clustering features & clusters & per-cluster & global & gain (as evaluated) & gain (matched) \\
\midrule
full history & 41 & 0.6937 & 0.6856 & 0.0081 & \textbf{0.0097} \\
first 75\% & 44 & 0.6942 & 0.6893 & 0.0049 & \textbf{0.0063} \\
first 50\% & 46 & 0.6937 & 0.6940 & -0.0003 & \textbf{0.0040} \\
first 25\% & 32 & 0.6777 & 0.6924 & -0.0147 & \textbf{-0.0057} \\
\bottomrule
\end{tabular}
\par\vspace{2pt}\footnotesize\raggedright Each row clusters on features truncated at a fraction of every case's own timeline, then predicts on the full held-out traces. The unmatched rows evaluate different populations, so their gradient mostly reflects the global baseline moving; the matched block rescores every truncation on the cases all four evaluate. The final column rescores every truncation on the 1,091 cases all four evaluate, which is the comparison the unmatched column cannot make.
\end{table}


%% GENERATED by `pipeline run results` - do not edit by hand.
\begin{table}[htbp]
\centering
\caption{Feature-block and temporal-dependency ablations.}\label{tab:ablation}
\begin{tabular}{lrrr}
\toprule
variant & ARI & NMI & clusters \\
\midrule
History truncated at 75\% & 0.798 & 0.895 & 68 \\
History truncated at 50\% & 0.649 & 0.816 & 70 \\
History truncated at 25\% & 0.467 & 0.717 & 53 \\
Keyword block removed & 0.596 & 0.838 & 65 \\
Topic block removed & 0.840 & 0.905 & 71 \\
\bottomrule
\end{tabular}
\par\vspace{2pt}\footnotesize\raggedright Each variant rebuilds the features under one change and re-clusters; agreement is measured against the full-history partition (64 clusters over 5,243 cases).
\end{table}


\subsection{What separates one partition from another, and what does not (RQ2)}\label{sec:res-rq2}

\subsubsection{Three distinct discards, and why raw comparison is invalid}

Coverage is lost three times, and conflating them produces the wrong conclusion. Density-based
clustering leaves 739 to 879 cases unassigned; centroid, modularity and consensus clustering
assign every case. Separately, the requirement of at least 30 traces per cluster removes
clusters from evaluation for every algorithm equally. And on the routing alphabet only, a case
whose itinerary contains fewer than two organisational units carries no transition to predict
and is dropped before either of the other two applies: \textbf{1{,}082 of the 5{,}243 cases,
20.6\%}, never enter that arm at all.

That third discard is alphabet-specific and it is the reason the two arms are not scored on the
same population --- the action alphabet is evaluated over all 5{,}243 cases, the routing
alphabet over 4{,}161. It also reapportions the loss the floor appears to cause. On the routing
alphabet, of the 1{,}187 to 2{,}207 cases absent from evaluation, roughly 1{,}082 are
single-unit cases rather than casualties of the floor, which leaves 58--77\% of the corpus
evaluated; consensus discards nothing at clustering time yet evaluates only about 60\%, because
fragmenting into 120 clusters puts many of them below the floor, while centroid clustering
reaches 77\%. On the action alphabet the floor removes only 105 to 1{,}156 cases and coverage
runs 77--98\%, so for most action-alphabet cells the floor costs \emph{less} than density
clustering's noise rather than more.

Because density-based clustering is scored only on the cases it chose to keep, comparing raw
per-algorithm numbers compares different populations. Rescoring every clustering on the 824
cases that all sixteen evaluate removes the confound without re-running anything, since
per-case correctness is stored alongside the held-out identifiers.

One property of that rescoring must be stated, because Section~\ref{sec:res-rq3} makes a great
deal of the same issue. The per-case record used here is the one belonging to whichever of the
five model families scored highest in that cluster, selected on the held-out data it is then
scored on. Every common-set figure in this subsection is therefore a per-cluster best-of-five,
and is an upper bound in the same sense as the selected-best figures of
Section~\ref{sec:res-rq3}. It is applied identically to all sixteen clusterings, so it does not
bias the comparison \emph{between} them, which is what this subsection uses it for; it would
bias any reading of the absolute level, and we do not offer one.

The apparent advantage disappears. On its own evaluated set the raw-arm density clustering
scores 0.6891, the best of the sixteen; on the common set it scores 0.6908, second-worst, with
only the coarse-parameter density clustering of the same arm below it. Every other clustering
gains between 0.020 and 0.033 under the restriction, because the restriction removes the
difficult cases that clustering had already discarded, while it gains 0.002 --- its own set was
already the easier subset.

The mechanism is narrower than ``density-based clustering'' as a family. Of the four density
clusterings, only the raw-arm one at 16.8\% noise shows the near-zero gain; the other three gain
0.020 to 0.023, indistinguishable from the partitional methods. What produces the reversal is a
high discard rate, not the algorithm family. Across all sixteen the spread on the common set is \textbf{0.008}.

Two qualifications bound what that number can be made to say. First, the comparison is made on
the cases every algorithm can model, which by the analysis in the next paragraph are the easier
ones --- restricting to the intersection removes precisely the cases on which the algorithms
might have differed, so this is a statement about the shared population, not about the corpus.

Second, and more restrictive, a small spread is not by itself evidence of equivalence, and this
study cannot convert it into one. A case-level bootstrap over the 824 common cases, 2{,}000
replicates, puts the spread at 0.0079 with a 95\% interval of $[0.0064, 0.0139]$ and a one-sided
90\% upper bound of 0.0124. An equivalence claim requires the bound, not the point estimate, to
fall inside a margin. Against the only margin this study registered --- 0.01 absolute top-1, and
registered for the anonymisation comparison on the action alphabet rather than for a sixteen-way
range on the routing one --- the bound does not fall inside it. We therefore report the spread
descriptively and make no equivalence claim: on the cases all algorithms can model, the observed
differences between clusterings are small, and this design cannot establish that they are
immaterial. Doing so would need a margin specified for this comparison and enough common cases
to tighten the interval, which is the same standard we hold the anonymisation result to in
Section~\ref{sec:res-rq5}.

That the discarded cases are harder, rather than unmodellable, can be checked directly. The
algorithms that assign every case score 0.556--0.585 on the cases density clustering rejects,
against 0.675--0.692 on the rest: a gap of roughly 0.10 to 0.14, or 14--20\% relatively worse.
For an artefact intended to describe an administration's actual work, silently omitting 14--17\%
of it is a real cost that buys no accuracy, and the discard should be read as a limitation of
that algorithm rather than as evidence of its quality.

%% GENERATED by `pipeline run results` - do not edit by hand.
\begin{table}[htbp]
\centering
\caption{Clustering comparison on a common set of cases.}\label{tab:common-set}
\footnotesize
\begin{tabular}{lrrrr}
\toprule
clustering & own $n$ & own set & common set & $\Delta$ \\
\midrule
exp-raw/leiden & 1,153 & 0.7107 & \textbf{0.7129} & 0.0022 \\
exp-raw-ksil/kmeans & 1,195 & 0.7081 & \textbf{0.7127} & 0.0046 \\
exp-anon/leiden & 1,122 & 0.7096 & \textbf{0.7114} & 0.0018 \\
exp-raw-kelbow/kmeans & 1,169 & 0.7126 & \textbf{0.7111} & -0.0015 \\
exp-raw-hcoarse/hdbscan & 1,207 & 0.7070 & \textbf{0.7109} & 0.0039 \\
exp-anon/hdbscan & 1,183 & 0.7054 & \textbf{0.7107} & 0.0053 \\
exp-anon-hcoarse/hdbscan & 1,199 & 0.7072 & \textbf{0.7101} & 0.0029 \\
exp-raw/hdbscan & 1,185 & 0.7030 & \textbf{0.7098} & 0.0068 \\
exp-anon/kmeans & 1,186 & 0.7058 & \textbf{0.7094} & 0.0036 \\
exp-raw-lres/leiden & 1,117 & 0.7046 & \textbf{0.7093} & 0.0047 \\
exp-anon-lres/leiden & 1,025 & 0.7049 & \textbf{0.7091} & 0.0042 \\
exp-raw/kmeans & 1,183 & 0.7067 & \textbf{0.7088} & 0.0021 \\
exp-anon-kelbow/kmeans & 1,170 & 0.7042 & \textbf{0.7074} & 0.0032 \\
exp-anon-ksil/kmeans & 1,195 & 0.7038 & \textbf{0.7074} & 0.0036 \\
\bottomrule
\end{tabular}
\par\vspace{2pt}\footnotesize\raggedright All clusterings rescored on the 967 cases every one of them evaluates. Spread on the common set: \textbf{0.0055}.
\end{table}


\subsubsection{Internal validity indices}

Table~\ref{tab:intrinsic} reports each index against downstream quality across the sixteen
clusterings. Both results in this subsection --- the index correlations and the common-set
comparison above --- are computed on the routing alphabet, and we state them as such rather than
as general claims about clustering. All sixteen indices are computed in the same space, the
reduced embedding in which clustering is performed. This is not automatic: the two consensus
partitions are derived from the other algorithms' label assignments rather than from the feature
matrix, and indices computed in that co-association representation are not on the same scale as
the rest --- a one-hot encoding of agreement makes members of a consensus cluster nearly
coincident, which inflates every compactness measure. We therefore recompute the consensus
partitions' indices in the shared embedding, and the evidence pipeline refuses to pool indices
carrying different space stamps.

Taken at face value the correlations are moderate and point the way the indices intend:
$\rho = 0.36$ for DBCV, $0.35$ for silhouette, $-0.48$ for Davies--Bouldin. None is significant
--- a permutation test over 20{,}000 draws rejects nothing, the smallest $p$ being 0.059 --- and
at sixteen observations the design rejects only for $|\rho| \geq 0.497$, so associations of the
size observed are exactly the ones it cannot resolve. Read no further and the result is an
underpowered null.

Reading further is what the data reward, because the association that does exist is not with
utility. Table~\ref{tab:index-mechanism} traces it. Across all sixteen clusterings, silhouette
is strongly related to \emph{how much a partition throws away}: $\rho = 0.61$ with the share of
cases discarded as noise. Discarding raises the score on the partition's own evaluated set
($\rho = 0.68$), because the cases a density method declines to assign are the hard ones
(Section~\ref{sec:res-rq2}). An index that rewards discarding, scored against a metric computed
on what survives the discarding, will correlate with it for reasons that have nothing to do with
cluster quality.

Restricting to the twelve clusterings that assign every case closes that channel by
construction, and one relationship in that subset is strong enough to survive the scrutiny a
sample of twelve deserves. Silhouette tracks granularity almost deterministically:
$\rho = 0.87$ with median cluster size and $-0.84$ with the number of clusters, and the
association holds on every subset we can form --- dropping the consensus partitions gives 0.92,
dropping the modularity ones 0.93, and the six centroid clusterings alone 0.83. On this corpus,
what silhouette measures is how coarse the partition is.

That finding is conditional on how the twelve were produced, and the condition is not incidental.
They are six centroid clusterings at different $k$, four modularity clusterings at different
resolutions, and two consensus partitions at a fixed $k$ --- so granularity is the variable the
parameter sweep deliberately manipulates. Finding that a compactness index tracks the knob the
sweep turns is closer to a property of that design than to a discovery about the corpus, and on
the six configuration-level points the association is $\rho = 0.77$ at $p = 0.10$ under an exact
permutation test over all orderings, which does not clear significance. What we can say is that on partitions generated this
way the index is a granularity measure; whether it would separate partitions of equal
granularity is a question this design cannot pose.

Two things that look like further findings are not, and we set them aside rather than count
them. Silhouette also correlates with coverage at $\rho = 0.84$, but coverage is determined by
how many clusters clear a size floor, so this is a restatement of the granularity result rather
than an addition to it: controlling for median cluster size, the partial correlation between
silhouette and coverage is 0.125. And silhouette correlates with accuracy at
$\rho = 0.35$ against $\rho = 0.84$ for coverage, which invites the reading that it predicts one
and not the other --- but twelve clusterings are six raw/anonymised pairs, and a bootstrap that
resamples configurations rather than clusterings puts the difference between those two
correlations at $[-0.02, +0.99]$ over 5{,}000 draws. The interval includes zero. We cannot distinguish them.

The same discipline applies to the null. At twelve observations the design resolves only
$|\rho| \geq 0.58$, so $\rho = 0.35$ is not evidence of a weak association any more than it is
evidence of a strong one; it is the range this design cannot see into. Stratifying by cluster
size leaves the parent partition's silhouette with no association with accuracy in any band
($-0.10$, $-0.01$, $-0.12$), but silhouette is a partition-level constant, so those bands
contain twelve distinct values rather than 488 independent ones, and the bands do not fully hold
size fixed --- cluster size still predicts accuracy within the smallest and largest of them.

What RQ2 supports is therefore one positive finding and one unchanged negative. The positive
finding is that on this corpus these indices are substantially a measure of granularity, which
is a real property and not the one they are usually invoked for. The negative finding is the
same as before: no evidence that they predict downstream utility, at a design that could not
have detected a moderate association had one existed. We do not claim to have shown they fail to
predict it. The practical consequence follows from the positive half alone --- selecting a
clustering on silhouette here selects, in substance, for coarseness, and
Section~\ref{sec:selection} therefore selects on the quantities directly.

%% GENERATED by `pipeline run results` - do not edit by hand.
\begin{table}[htbp]
\centering
\caption{What the validity indices actually track.}\label{tab:index-mechanism}
\footnotesize
\begin{tabular}{lll}
\toprule
relation & population & $\rho$ \\
\midrule
Silhouette vs.\ share of cases discarded as noise & all 14 & -0.308 \\
Share discarded vs.\ accuracy on the partition's own set & all 14 & -0.204 \\
Silhouette vs.\ median cluster size & full-cov.\ 14 & 0.427 \\
Silhouette vs.\ number of clusters & full-cov.\ 14 & -0.325 \\
Number of clusters vs.\ coverage & full-cov.\ 14 & -0.991 \\
Silhouette vs.\ coverage & full-cov.\ 14 & 0.341 \\
\textbf{Silhouette vs.\ common-set accuracy} & full-cov.\ 14 & \textbf{-0.392} \\
\textbf{DBCV vs.\ common-set accuracy} & full-cov.\ 14 & \textbf{-0.057} \\
Silhouette vs.\ accuracy, small clusters only (mean top-1 0.623) & 233 clusters, 14 values & -0.151 \\
Silhouette vs.\ accuracy, medium clusters only (mean top-1 0.655) & 267 clusters, 14 values & -0.009 \\
Silhouette vs.\ accuracy, large clusters only (mean top-1 0.781) & 89 clusters, 14 values & 0.074 \\
\bottomrule
\end{tabular}
\par\vspace{2pt}\footnotesize\raggedright Granularity is the relationship that survives scrutiny at this sample size. The coverage rows are largely a restatement of it (partial correlation given median cluster size is small), and the difference between the coverage and accuracy correlations is not resolvable here -- its bootstrap interval over configurations includes zero. The band rows carry only as many distinct silhouette values as there are partitions, not as many as there are clusters, so they are far weaker evidence than their cluster counts suggest.
\end{table}


This result concerns indices computed on a reduced embedding, which is where clustering is
performed; it is not a claim about the embedding itself. It also does not contradict work
certifying these indices against recovery of a known partition, since no known partition exists
here.

%% GENERATED by `pipeline run results` - do not edit by hand.
\begin{table}[htbp]
\centering
\caption{Internal validity indices against downstream utility.}\label{tab:intrinsic}
\begin{tabular}{lrr}
\toprule
index & $\rho$ & $p$ \\
\midrule
DBCV & 0.538 & 0.053 \\
Silhouette & -0.516 & 0.060 \\
Davies--Bouldin & 0.464 & 0.100 \\
Calinski--Harabasz & -0.270 & 0.354 \\
Noise ratio (refused) & -0.055 & 0.853 \\
Number of clusters & 0.501 & 0.073 \\
Median cluster size & -0.524 & 0.061 \\
\bottomrule
\end{tabular}
\par\vspace{2pt}\footnotesize\raggedright Spearman rank correlation with downstream quality across 14 clusterings; permutation test with 20,000 draws. The downstream measure spans 0.112 (34\% relative).
\end{table}


\subsection{Typical paths against three controls (RQ3)}\label{sec:res-rq3}

\subsubsection{The partition carries the signal}

Permuting cluster membership while preserving cluster count, size distribution and evaluated
population reduces best next-step accuracy from 0.726 to 0.494 on the action alphabet, and from
0.672 to 0.457 on the routing alphabet. Nothing about the \emph{shape} of a partition produces
23 points; only membership does. The same conclusion holds without any learned model: the
within-cluster most-frequent-successor lookup scores 0.643 on real clusters and 0.447 on
shuffled ones, which is evidence of coherence independent of the sequence models entirely.

\subsubsection{What per-cluster modelling adds, measured honestly}

Table~\ref{tab:controls} reports all three baselines, and the comparison must be made carefully,
because the obvious way to make it is wrong.

Taking the best of the five model families in each cluster gives 0.726 on actions and 0.672 on
routing, against 0.713 and 0.610 for a single global third-order Markov chain: apparent gains of
0.013 and 0.063. That comparison is not like-for-like. The clustered figure is a per-cluster
\emph{maximum} over five families, selected using held-out data, while the global baseline is
one fixed family. Part of the apparent gain is the selection, not the partitioning.

Matching the families removes it. A third-order Markov chain fitted per cluster, against the
same chain fitted globally, gains \textbf{0.004} on actions and \textbf{0.035} on routing.
Fixing any single family throughout rather than choosing per cluster gives at most 0.005 on
actions and 0.053 on routing. On actions a per-cluster LSTM is \emph{worse} than the global
model \((-0.004)\).

We report both: the selected-best figures because they are what a practitioner choosing freely
would obtain, and the family-matched figures because they are what the partition itself
contributes. The honest statement of the second is that \textbf{on the action alphabet,
partitioning contributes almost nothing} --- 0.004 against a baseline of 0.713.

This exposes a symmetry worth naming. The paper's own argument is that the conventional
within-cluster baseline inflates the apparent contribution of clustering roughly sixfold;
selecting the best of five families against a single-family baseline inflates it again, roughly
threefold on actions. The same discipline has to apply at both levels.

The gap between the two alphabets is the substantive finding, and it must be measured on one
population before it can be stated. As reported so far the two arms are scored on different
populations: the 1{,}082 single-unit cases that carry no itinerary are absent from the routing
arm by construction (Section~\ref{sec:res-rq2}). Those are the simplest cases in the corpus, so
excluding them from routing while retaining them in actions biases the comparison in the
direction that flatters routing, and a ratio computed across the two is not interpretable.

We therefore re-fitted the action-alphabet models on the 4{,}161 multi-unit cases, so that both
alphabets are scored over the identical 323 clusters and 27{,}795 traces.
Table~\ref{tab:shared-population} reports the result, and the correction does not weaken the
finding --- it inverts one half of it. On all cases, a per-cluster Markov chain beats a global
one on actions by 0.0032 over these eight clusterings, the same near-zero gain the full sixteen
put at 0.004. On the shared population it \emph{loses} to the global model by 0.0087, and it
does so in all eight clusterings individually, from $-0.0020$ to $-0.0133$. Routing on that same
population gains 0.0359.

The mechanism is visible in the same table. Removing the single-unit cases costs the global model
0.011 of accuracy and the per-cluster models 0.023 --- twice as much. The small positive gain
that per-cluster action models showed was drawn disproportionately from the easiest cases in the
corpus; on the cases that actually have a routing decision to make, partitioning does not help
action prediction, it hurts it.

The substantive claim is therefore signed rather than scaled: \textbf{on one population,
per-cluster modelling helps on routing and does not help on actions}. Both magnitudes depend on
a choice we must expose rather than bury, and we expose it for both rather than only where it
is convenient.

Aggregating per-cluster accuracies to a corpus figure requires a weight, and the two defensible
weights disagree. Weighting each cluster by its membership --- which asks how a randomly chosen
\emph{case} is served --- gives $+0.0032$ on actions over all cases, $-0.0087$ on actions over
the shared population, and $+0.0359$ on routing. Weighting by held-out prediction steps --- the
pooled accuracy over every prediction actually made --- gives $-0.0044$, $-0.0047$ and
$+0.0292$. Two things follow. On actions the second weighting removes the sign change entirely:
per-cluster modelling is slightly negative on both populations, and restricting the population
moves it by 0.0003 rather than by 0.012, so the apparent inversion is an artefact of the weight
and not a population effect. On routing the gain is real under both weights but 18\% smaller
under the second, so the headline figure is weight-dependent in the same way and to a
comparable degree.

What survives either choice is the contrast and its direction: routing gains something between
0.029 and 0.036, actions gain nothing and may lose a little. That is the finding; the third
decimal place is not.

Re-fitting every family on the restricted population shows this is not a property of the Markov
chain (Table~\ref{tab:families-shared}). Membership-weighted, all five are negative there. The
selected-best figure behaves the same way: the per-cluster maximum over five families gains
0.0116 over a global model on all cases and \textbf{0.0007} on the shared population --- and
that maximum is an oracle, chosen on the same held-out data it is scored on, so it is an upper
bound rather than an estimate. An oracle over five families barely reaching parity with one
pooled chain is the clearest statement of the result.

We do not report an interval on the action deficit, and it is small enough that one is needed
before any strong claim rests on it; what the evidence supports is a sign and a recommendation,
not a magnitude. Which unit handles a case next depends on the procedure; which action is recorded
largely follows conventions shared across procedures. For an instrument whose purpose is to show
how work moves through an administration, the routing representation is both the more
informative and the better supported.

Size changes the picture on actions and not on routing. On actions the selected-best gain over a
global model is 0.003, 0.010 and 0.024 for small, medium and large clusters; since the pooled
selected-best gain of 0.013 corresponds to a family-matched 0.004, the family-matched gain below
150 traces is indistinguishable from zero, and a per-cluster action model is not worth training
there. On routing it is 0.055, 0.072 and 0.052 ---
essentially size-independent, and worth training throughout. Fabrication follows the same
split: unseen transitions fall from 6.04 to 1.10 with size on actions, but are already low on
routing and reach 0.038 on large clusters, where the generated path invents essentially nothing.

%% GENERATED by `pipeline run results` - do not edit by hand.
\begin{table}[htbp]
\centering
\caption{Next-step accuracy against three baselines, action-by-unit alphabet.}\label{tab:controls-unit}
\begin{tabular}{rrrrrrrrr}
\toprule
sel.\ best & successor & global & shuffled & $\Delta$ fixed family & $\Delta$ best family & edge-F1 & edge-F1 shuffled & $\Delta^{*}$ edge-F1 \\
\midrule
0.722 & 0.632 & 0.712 & 0.512 & \textbf{0.000} & 0.003 & 0.501 & 0.361 & 0.140 \\
\bottomrule
\end{tabular}
\par\vspace{2pt}\footnotesize\raggedright ``selected best'' takes the per-cluster maximum over every model family, so its gain over a single-family global baseline is not like-for-like; the $\Delta$ columns hold the family fixed. $\Delta$ still contains whatever a partition of this size distribution earns on its own, which is why the shuffled column is reported beside it and the partition's own contribution is separated into its own table. edge-F1 is the harmonic mean of directly-follows edge coverage and edge usage; it is a cluster statistic rather than a model output, so it takes one value per partition and cannot vary by model family. Every number is measured inside one alphabet, against a baseline facing the same vocabulary; alphabets are never tabulated together, because accuracies over vocabularies of different size are not the same quantity.
\end{table}

\begin{table}[htbp]
\centering
\caption{Next-step accuracy against three baselines, routing alphabet.}\label{tab:controls-itinerary}
\begin{tabular}{rrrrrrrrr}
\toprule
sel.\ best & successor & global & shuffled & $\Delta$ fixed family & $\Delta$ best family & edge-F1 & edge-F1 shuffled & $\Delta^{*}$ edge-F1 \\
\midrule
0.721 & 0.481 & 0.689 & 0.517 & \textbf{0.010} & 0.023 & 0.545 & 0.425 & 0.120 \\
\bottomrule
\end{tabular}
\par\vspace{2pt}\footnotesize\raggedright ``selected best'' takes the per-cluster maximum over every model family, so its gain over a single-family global baseline is not like-for-like; the $\Delta$ columns hold the family fixed. $\Delta$ still contains whatever a partition of this size distribution earns on its own, which is why the shuffled column is reported beside it and the partition's own contribution is separated into its own table. edge-F1 is the harmonic mean of directly-follows edge coverage and edge usage; it is a cluster statistic rather than a model output, so it takes one value per partition and cannot vary by model family. Every number is measured inside one alphabet, against a baseline facing the same vocabulary; alphabets are never tabulated together, because accuracies over vocabularies of different size are not the same quantity.
\end{table}

\begin{table}[htbp]
\centering
\caption{Next-step accuracy against three baselines, action alphabet.}\label{tab:controls-procedure}
\begin{tabular}{rrrrrrrrr}
\toprule
sel.\ best & successor & global & shuffled & $\Delta$ fixed family & $\Delta$ best family & edge-F1 & edge-F1 shuffled & $\Delta^{*}$ edge-F1 \\
\midrule
0.760 & 0.649 & 0.624 & 0.607 & \textbf{0.109} & 0.131 & 0.729 & 0.734 & -0.006 \\
\bottomrule
\end{tabular}
\par\vspace{2pt}\footnotesize\raggedright ``selected best'' takes the per-cluster maximum over every model family, so its gain over a single-family global baseline is not like-for-like; the $\Delta$ columns hold the family fixed. $\Delta$ still contains whatever a partition of this size distribution earns on its own, which is why the shuffled column is reported beside it and the partition's own contribution is separated into its own table. edge-F1 is the harmonic mean of directly-follows edge coverage and edge usage; it is a cluster statistic rather than a model output, so it takes one value per partition and cannot vary by model family. Every number is measured inside one alphabet, against a baseline facing the same vocabulary; alphabets are never tabulated together, because accuracies over vocabularies of different size are not the same quantity.
\end{table}


%% GENERATED by `pipeline run results` - do not edit by hand.
\begin{table}[htbp]
\centering
\caption{Family-matched Markov against a global model, on one population.}\label{tab:shared-population}
\begin{tabular}{lrrrrr}
\toprule
comparison & clusters & traces & per-cluster & global & gain \\
\midrule
Action, all cases & 429 & 36,051 & 0.7158 & 0.7126 & \textbf{0.0032} \\
Action, multi-unit cases only & 323 & 27,795 & 0.6928 & 0.7015 & \textbf{-0.0087} \\
Routing (multi-unit by construction) & 323 & 27,795 & 0.6448 & 0.6089 & \textbf{0.0359} \\
\bottomrule
\end{tabular}
\par\vspace{2pt}\footnotesize\raggedright Family-matched third-order Markov, fitted per cluster and scored against a single global chain, over the eight main clusterings. The routing alphabet excludes single-unit cases by construction; the second row applies that same exclusion to the action alphabet, so the last two rows are scored on one population of 4,161 cases.
\end{table}


%% GENERATED by `pipeline run results` - do not edit by hand.
\begin{table}[htbp]
\centering
\caption{Per-cluster gain by model family, before and after equalising the population.}\label{tab:families-shared}
\begin{tabular}{lrr}
\toprule
family & all cases & multi-unit only \\
\midrule
Markov-3 & 0.0030 & \textbf{-0.0087} \\
GRU & 0.0022 & \textbf{-0.0083} \\
LSTM & -0.0071 & \textbf{-0.0168} \\
Transformer & 0.0026 & \textbf{-0.0083} \\
Directly-follows & -0.0687 & \textbf{-0.0957} \\
\emph{best of five, selected per cluster} & 0.0116 & \textbf{0.0007} \\
\bottomrule
\end{tabular}
\par\vspace{2pt}\footnotesize\raggedright Gain over a single global third-order Markov chain on the action alphabet, by model family. The right column restricts to the cases the routing alphabet can also be scored on. Every family is positive-to-negligible on the left and negative on the right; the selected-best row is a per-cluster maximum over five families chosen on the same held-out data it is scored on, so it is an upper bound rather than an estimate.
\end{table}


\subsubsection{Against a real case}

The medoid is a strong baseline, and we report it on the same footing as the models rather than
as a foil. Its construction should be stated precisely, because it bears on how strong: it is
the trace with the smallest mean normalised edit distance to a random sample of at most twenty
traces drawn from the cluster's training split, not to the cluster as a whole. For clusters well
above that size the medoid is therefore a weaker representative than the name implies, and the
comparisons below understate it. On the five support measures of Section~\ref{sec:protocol} it
wins three on the action alphabet; on routing it retains the best unit recall (0.763 against
0.740 for the best model).

The measure on which the generated paths win is the one that matters for the claim, and it wins
on one alphabet only. Counting transitions that occur in \emph{no} sampled member of the cluster
--- fabrication rather than variation --- the third-order Markov path scores 0.227 on routing
against the medoid's 0.352. Normalising by path length, which prevents a short path from
appearing faithful merely by generating less, gives 0.0194 unseen transitions per visit for
Markov against 0.0252 for the medoid. A generated \emph{routing} path is less idiosyncratic than
any single real case.

The same normalisation reverses on the action alphabet, and reporting it only where it favours
the models would be the selective reporting this paper argues against elsewhere. There the
medoid fabricates 0.187 unseen transitions per visit against 0.327 for the best model
(Transformer), 0.351 for the LSTM and 0.630 for the Markov chain --- the medoid is better by a
factor of 1.8 to 3.4, on the same measure and the same normalisation. The claim that a generated
path is less idiosyncratic than a real one therefore holds for the routing representation and
not for the action one. That is the narrower claim the evidence supports, and it is the
representation the rest of the paper argues matters for transparency; but the qualifier is
required, and the medoid remains the better representative where actions are concerned.

Model behaviour follows their capacity. The Markov chain is the most conservative generator; the
Transformer produces the longest paths (12.6 visits against the medoid's 13.9) with the best
unit recall and nearly the lowest per-visit fabrication. On large clusters the neural models
overtake Markov, which its third-order memory cannot represent; on small clusters they overfit
and fabricate more, and the short memory acts as a regulariser. Every model under-generates:
matched visit fractions cluster around 0.45, meaning fewer than half of a real case's visits are
covered by any generated path. This is the clearest open weakness in the task.

%% GENERATED by `pipeline run results` - do not edit by hand.
\begin{table}[htbp]
\centering
\caption{Path quality by cluster size, action-by-unit alphabet.}\label{tab:size-trend-unit}
\begin{tabular}{lrrrrr}
\toprule
cases in cluster & clusters & top-1 & gain vs global & unit recall & unseen \\
\midrule
30-50 & 183 & 0.669 & 0.0047 & 0.705 & 8.154 \\
50-75 & 232 & 0.693 & 0.0062 & 0.643 & 2.581 \\
75-100 & 130 & 0.669 & -0.0015 & 0.650 & 3.401 \\
100-150 & 109 & 0.725 & 0.0170 & 0.711 & 3.273 \\
150-250 & 85 & 0.762 & 0.0108 & 0.774 & 1.331 \\
250- & 25 & 0.817 & 0.0283 & 0.795 & 0.252 \\
\bottomrule
\end{tabular}
\par\vspace{2pt}\footnotesize\raggedright Gain is measured band by band against a global model of the same family. Clusters below thirty traces are not evaluated.
\end{table}

\begin{table}[htbp]
\centering
\caption{Path quality by cluster size, routing alphabet.}\label{tab:size-trend-itinerary}
\begin{tabular}{lrrrrr}
\toprule
cases in cluster & clusters & top-1 & gain vs global & unit recall & unseen \\
\midrule
30-50 & 157 & 0.668 & 0.0361 & 0.737 & 1.226 \\
50-75 & 148 & 0.657 & 0.0264 & 0.648 & 0.554 \\
75-100 & 111 & 0.665 & 0.0543 & 0.670 & 0.470 \\
100-150 & 84 & 0.719 & 0.0383 & 0.702 & 0.393 \\
150-250 & 70 & 0.782 & 0.0172 & 0.811 & 0.268 \\
250- & 19 & 0.830 & 0.0257 & 0.797 & 0.103 \\
\bottomrule
\end{tabular}
\par\vspace{2pt}\footnotesize\raggedright Gain is measured band by band against a global model of the same family. Clusters below thirty traces are not evaluated.
\end{table}

\begin{table}[htbp]
\centering
\caption{Path quality by cluster size, action alphabet.}\label{tab:size-trend-procedure}
\begin{tabular}{lrrrr}
\toprule
cases in cluster & clusters & top-1 & gain vs global & unseen \\
\midrule
30-50 & 183 & 0.721 & 0.1092 & 4.215 \\
50-75 & 232 & 0.743 & 0.1340 & 1.794 \\
75-100 & 130 & 0.718 & 0.1229 & 1.990 \\
100-150 & 109 & 0.757 & 0.1313 & 2.675 \\
150-250 & 85 & 0.793 & 0.1450 & 2.817 \\
250- & 25 & 0.827 & 0.1797 & 5.155 \\
\bottomrule
\end{tabular}
\par\vspace{2pt}\footnotesize\raggedright Gain is measured band by band against a global model of the same family. Clusters below thirty traces are not evaluated.
\end{table}


\subsection{Rule compliance and the value of a guarantee (RQ4)}\label{sec:res-rq4}

Unconstrained generation is already largely valid against the send/receive pairing rule that
governs case movement. Between 98.8\% and 99.8\% of generated steps satisfy it, and between
89.2\% and 96.3\% of whole generated traces do.

The comparison that gives this meaning is the event log itself, which satisfies the rule on
92.06\% of its traces. Roughly one real case in twelve breaks it. The generators fall on both
sides of that line and the split should be given in full rather than by its favourable half:
the three neural families are \emph{less} rule-compliant than observed practice (Transformer
89.40\%, LSTM 89.18\%, GRU 90.93\%), while the Markov chain at 95.65\% and the directly-follows
graph at 96.32\% are more so. On the routing alphabet the question does not arise: validity is
1.0 for every model and for the log, because that alphabet carries no pairing to violate.

The contribution of constrained decoding is therefore not rescuing models from widespread
invalidity. It is converting a high but uncertain compliance rate into a guarantee that holds by
construction for the property the automaton encodes, which is what an artefact intended for
audit requires: a catalogue entry that might violate the procedure's own rules is not auditable.
Two limits belong in the same breath. The constraint enforces a rule that observed practice
breaks about 8\% of the time, so a constrained path describes the procedure as it is meant to
run, not always as it does. And the guarantee is by construction rather than by measurement:
Table~\ref{tab:flow} reports unconstrained generation, because a constrained path satisfies the
automaton by definition and rescoring it would measure the decoder rather than the model. That
makes the claim analytic, and it therefore inherits exactly the bounds set out in
Section~\ref{sec:method} --- no unmatched receive and no unopened conclusion, over the five
action types the rules govern, with termination itself unconstrained.

What the constraint costs can be measured, and on the conservative generator the answer is that
it costs nothing because it does nothing. Walking the held-out traces of one clustering and
comparing, at each of 69{,}012 decision points, the third-order Markov chain's top-eight
candidates with and without the automaton's filter: the filter alters the candidate set at
\textbf{0.0\%} of steps, changes the most likely candidate at 0.0\%, and never leaves the model
with no admissible move. The unconstrained argmax is never illegal. The reason is structural
rather than fortunate --- a Markov chain's next-state distribution contains only transitions
observed in training, and observed transitions are overwhelmingly legal, so the model rarely has
an illegal continuation to propose in the first place.

This sharpens rather than undercuts the section's claim, and it also bounds it. The guarantee is
free for this generator, which is why constrained decoding can be adopted without an accuracy
argument. But a guarantee that never binds is also not what is rescuing the artefact, and the
invalidity the flow table does record --- 4\% to 11\% of whole generated traces --- arises when
generation runs free of a real prefix and compounds its own output, not from single illegal
steps. The measurement above covers the Markov generator at real prefixes on one clustering; the
neural families, whose whole-trace validity falls below the log's, are where the constraint
would be expected to bind, and we have not measured them.

%% GENERATED by `pipeline run results` - do not edit by hand.
\begin{table}[htbp]
\centering
\caption{Rule compliance of unconstrained generation, against the event log itself.}\label{tab:flow}
\begin{tabular}{lrr}
\toprule
source & valid steps & fully valid traces \\
\midrule
markov & 0.9978 & 0.9486 \\
dfg & 0.9934 & 0.9337 \\
gnn & 0.9869 & 0.8756 \\
gru & 0.9930 & 0.8894 \\
lstm & 0.9908 & 0.8639 \\
transformer & 0.9900 & 0.8719 \\
\emph{real event log} & 0.9986 & 0.9162 \\
\bottomrule
\end{tabular}
\par\vspace{2pt}\footnotesize\raggedright Constrained decoding is valid by construction and is not tabulated.
\end{table}


\subsection{Anonymisation (RQ5)}\label{sec:res-rq5}

Comparing the raw and anonymised arms case by case, over 1{,}143 to 1{,}515 cases held out in
both depending on the clustering, a
paired bootstrap with 2{,}000 replicates finds all 20 model-by-clustering comparisons
non-inferior at the pre-registered margin $\delta = 0.01$, registered before the data existed.

At the tighter sensitivity margin $\delta = 0.005$ the result is \textbf{17 of 20}. The three
failures are the neural models on the centroid clustering --- GRU, LSTM and Transformer, with
one-sided 90\% lower bounds of $-0.0059$, $-0.0063$ and $-0.0072$ against a threshold of
$-0.005$. They fail narrowly, and only under a margin half the pre-registered one, but they
fail; reporting the sensitivity analysis as though it had passed would be the selective
reporting this paper argues against elsewhere.

The claim this supports is narrow, and prior work sets its bounds. The utility cost of event-log
sanitisation for process discovery is established \citep{fahrenkrogpetersen2019pretsa}, and the
effect of anonymisation on clustering quality has been measured on administrative data
\citep{ferrao2022utility}. What we add is the effect on \emph{downstream next-activity
prediction}, tested against a pre-registered non-inferiority margin rather than reported as a
difference in indices. Within that scope, and at the pre-registered margin, anonymisation
applied before clustering costs no measurable utility; at half that margin the conclusion holds
for 17 of 20 comparisons.
